
\begin{nobreaks}
\subsubsection*{Matrices and piecewise functions}
\phantomsection\label{matrices}
\begin{referencecode}
\[
    A = \begin{bmatrix}
        1 & 2 \\
        3 & 4
    \end{bmatrix}
    \qquad
    \det(A) = \begin{vmatrix}
        1 & 2 \\
        3 & 4
    \end{vmatrix} = -2
\]
\end{referencecode}

\[
    A = \begin{bmatrix}
        1 & 2 \\
        3 & 4
    \end{bmatrix}
    \qquad
    \det(A) = \begin{vmatrix}
        1 & 2 \\
        3 & 4
    \end{vmatrix} = -2
\]

\texttt{pmatrix}: parentheses; \texttt{bmatrix}: brackets;
\texttt{vmatrix}: determinant bars; \texttt{matrix}: no delimiters.
\end{nobreaks}

\begin{referencecode}
\[
    f(x) = \begin{cases}
        x^2, & \text{if } x \geq 0, \\
        -x,  & \text{if } x < 0.
    \end{cases}
\]
\end{referencecode}

\[
    f(x) = \begin{cases}
        x^2, & \text{if } x \geq 0, \\
        -x,  & \text{if } x < 0.
    \end{cases}
\]

\begin{referencecode}
\[
    \left(\frac{a+b}{c}\right)^2
    \qquad \sin(\theta)
    \qquad \mathbf{v} = (v_x,v_y)
\]
\end{referencecode}

\[
    \left(\frac{a+b}{c}\right)^2
    \qquad \sin(\theta)
    \qquad \mathbf{v} = (v_x,v_y)
\]

Use \texttt{\string\text\{...\}} for words and commands such as
\texttt{\string\sin} or \texttt{\string\log} for operators.

\partlabel{Systems with taller rows}
The template extends \texttt{cases} with optional \texttt{[L]} and
\texttt{[XL]} row struts. Its definition is in
\path{Packages/Math/package_config.tex} and uses \texttt{array}
and the presets loaded by the Graphics package group.
\begin{referencecode}
\[
\begin{cases}[L]
    \zeta\omega_n = \dfrac{4-a}{2}, \\
    \omega_n^2 = \dfrac{K}{a}.
\end{cases}
\]
\end{referencecode}

\[
\begin{cases}[L]
    \zeta\omega_n = \dfrac{4-a}{2}, \\
    \omega_n^2 = \dfrac{K}{a}.
\end{cases}
\]

Omit the option for ordinary \texttt{cases}. Both forms use a local
\texttt{\string\arraystretch} of 1.2. The optional presets are
template extensions, not standard \texttt{amsmath} syntax.
\partlabel{A brace around equations aligned at the equals sign}
Unlike the left-aligned columns in \texttt{cases}, \texttt{aligned}
lets the alignment point sit at each equals sign. Place it inside
display math and close the left brace with the invisible delimiter
\verb|\right.|.
\begin{referencecode}
\[
\left\{
\begin{aligned}
    a^4 - 24a^2 + 15 + 2K &= 0 \\
    8a^3 - (32+K)a &= 0
\end{aligned}
\right.
\]
\end{referencecode}

\[
\left\{
\begin{aligned}
    a^4 - 24a^2 + 15 + 2K &= 0 \\
    8a^3 - (32+K)a &= 0
\end{aligned}
\right.
\]
